\par\Needspace{6\baselineskip}\medskip
\begin{center}\hypertarget{native-231-article168}{}\pdfbookmark[2]{168.}{article-168}168.\end{center}
Combining equations 1 and 2 of Article 112 gives the following equations in the notation used here.
\[
\begin{aligned}
0={}&(n12)R\cos B\sin(l-a)\\
&-(n02)\bigl(\rho'\cos\beta'\sin(a'-a)+R'\cos B'\sin(l'-a)\bigr)\\
&+(n01)\bigl(\rho''\cos\beta''\sin(a''-a)+R''\cos B''\sin(l''-a)\bigr)
\end{aligned}
\]
\SourceLeaf{232}
\[
\begin{aligned}
0={}&(n23)\bigl(\rho'\cos\beta'\sin(a'''-a')+R'\cos B'\sin(a'''-l')\bigr)\\
&-(n13)\bigl(\rho''\cos\beta''\sin(a'''-a'')+R''\cos B''\sin(a'''-l'')\bigr)\\
&+(n12)R'''\cos B'''\sin(a'''-l''').
\end{aligned}
\]

If we set
\[
\begin{gathered}
\frac{R'\cos B'\sin(l'-a)}{\cos\beta'\sin(a'-a)}-R'\cos\delta'=b'\\[.7em]
\frac{R\cos B\sin(l-a)}{\cos\beta''\sin(a''-a)}=\lambda\\[.7em]
\frac{R''\cos B''\sin(a'''-l'')}{\cos\beta''\sin(a'''-a'')}-R''\cos\delta''=b''\\[.7em]
\frac{R'''\cos B'''\sin(a'''-l''')}{\cos\beta'\sin(a'''-a')}=\lambda'''\\[.7em]
\frac{R'\cos B'\sin(a'''-l')}{\cos\beta'\sin(a'''-a')}-R'\cos\delta'=\varkappa'\\[.7em]
\frac{\cos\beta'\sin(a'-a)}{\cos\beta''\sin(a''-a)}=\mu'\\[.7em]
\frac{R''\cos B''\sin(l''-a)}{\cos\beta''\sin(a''-a)}-R''\cos\delta''=\varkappa''\\[.7em]
\frac{\cos\beta''\sin(a'''-a'')}{\cos\beta'\sin(a'''-a')}=\mu''
\end{gathered}
\]
and carry out the necessary reductions, these equations take the form
\[
\begin{gathered}
\frac{\mu'(1+P')(x'+b')}{1+\dfrac{Q'}{(x'x'+a'a')^{\frac32}}}=x''+\varkappa''+\lambda P'\\[1em]
\frac{\mu''(1+P'')(x''+b'')}{1+\dfrac{Q''}{(x''x''+a''a'')^{\frac32}}}=x'+\varkappa'+\lambda'''P'',
\end{gathered}
\]
or, if we further set
\[
\begin{gathered}
-\varkappa''-\lambda P'=c',\qquad\mu'(1+P')=d'\\[.5em]
-\varkappa'-\lambda'''P''=c'',\qquad\mu''(1+P'')=d'',
\end{gathered}
\]
they take the form
\[
\tag*{I.}x''=c'+\frac{d'(x'+b')}{1+\dfrac{Q'}{(x'x'+a'a')^{\frac32}}}
\]
\[
\tag*{II.}x'=c''+\frac{d''(x''+b'')}{1+\dfrac{Q''}{(x''x''+a''a'')^{\frac32}}}\mathpunct{\cdot}
\]
These two equations enable us to determine $x'$ and $x''$ from \newline $a',b',c',d',Q',a'',b'',c'',d'',Q''$. If we had to eliminate $x'$ or $x''$ from them, we would arrive at an equation of very high degree\HistoricalNoteEight{Manuscript note:] of degree $64^{\mathrm{ti}}$.}. Indirect methods, however, yield the values of the unknowns $x',x''$ fairly \SourceLeaf{233}readily without changing the form of these equations. In most cases we can already obtain approximate values of the unknowns by initially neglecting $Q'$ and $Q''$. This gives
\[
\begin{gathered}
x'=\frac{c''+d''(b''+c')+d'd''b'}{1-d'd''}\\[.6em]
x''=\frac{c'+d'(b'+c'')+d'd''b''}{1-d'd''}\mathpunct{\cdot}
\end{gathered}
\]
As soon as an approximate value of either unknown is available, values that satisfy the equations exactly can be found very easily. Let $\xi'$ be an approximate value of $x'$ that gives $x''=\xi''$ when substituted in equation I. Likewise, let substituting $x''=\xi''$ in equation II give $x'=X'$. Repeat the same operations with $x'$ in I set to a different value $\xi'+\nu'$ to obtain $x''=\xi''+\nu''$. Substituting this value in II gives $x'=X'+N'$. Then the corrected value of $x'$ is
\[
=\xi'+\frac{(\xi'-X')\nu'}{N'-\nu'}=\frac{\xi'N'-X'\nu'}{N'-\nu'}\mathpunct{\raisebox{.45ex}{,}}
\]
and the corrected value of $x''$
\[
=\xi''+\frac{(\xi'-X')\nu''}{N'-\nu'}\mathpunct{\cdot}
\]
If it seems worthwhile, the same operations are repeated with the corrected value of $x'$ and another value differing slightly from it, until values of $x',x''$ are obtained that satisfy equations I and II exactly. In any case, anyone with even moderate practice in analysis will find ways to shorten the calculation.

In these operations the irrational quantities $(x'x'+a'a')^{\frac32}$, $(x''x''+a''a'')^{\frac32}$ can be calculated conveniently by introducing the arcs $z',z''$, whose tangents are respectively $\dfrac{a'}{x'}\mathpunct{\raisebox{.45ex}{,}}$ $\dfrac{a''}{x''}\mathpunct{\raisebox{.45ex}{,}}$ whence
\[
\begin{gathered}
\surd(x'x'+a'a'\mathpunct{\raisebox{-.5ex}{.}})=r'=\frac{a'}{\sin z'}=\frac{x'}{\cos z'}\\[.6em]
\surd(x''x''+a''a'')=r''=\frac{a''}{\sin z''}=\frac{x''}{\cos z''}\mathpunct{\cdot}
\end{gathered}
\]
These auxiliary arcs, which must be taken between $0$ and $180^0$ so that $r',r''$ are positive, are evidently identical with the arcs $C'B'$, $C''B''$. Thus this procedure gives $r'$ and $r''$ as well as the positions of the points $C',C''$.

This determination of the quantities $x',x''$ requires that\newline $a',a'',b',b'',c',c'',d',d'',Q',Q''$ be known. The first four of these quantities are obtained from the data of the problem, while the next four depend on $P',P''$. The \SourceLeaf{234}quantities $P',P'',Q',Q''$ cannot yet be determined exactly. We do, however, have
\[
\tag*{III.}P'=\frac{t''-t'}{t'-t}\cdot\frac{(\eta01)}{(\eta12)}
\]
\[
\tag*{IV.}P''=\frac{t''-t'}{t'''-t''}\cdot\frac{(\eta23)}{(\eta12)}
\]
\[
\tag*{V.}\begin{aligned}
Q'={}&\tfrac12kk(t'-t)(t''-t')\cdot\frac{r'r'}{rr''}\cdot{}\\[.5em]
&\frac{1}{(\eta01)(\eta12)\cos\tfrac12(v'-v)\cos\tfrac12(v''-v)\cos\tfrac12(v''-v')}
\end{aligned}
\]
\[
\tag*{VI.}\begin{aligned}
Q''={}&\tfrac12kk(t''-t')(t'''-t'')\cdot\frac{r''r''}{r'r'''}\cdot{}\\[.5em]
&\frac{1}{(\eta12)(\eta23)\cos\tfrac12(v''-v')\cos\tfrac12(v'''-v')\cos\tfrac12(v'''-v'')}\mathpunct{\raisebox{.45ex}{,}}
\end{aligned}
\]
so we immediately have the approximate values
\[
\begin{gathered}
P'=\frac{t''-t'}{t'-t}\mathpunct{\raisebox{.45ex}{,}}\qquad P''=\frac{t''-t'}{t'''-t''}\\[.6em]
Q'=\tfrac12kk(t'-t)(t''-t'),\qquad Q''=\tfrac12kk(t''-t')(t'''-t''),
\end{gathered}
\]
on which the first calculation will be based.
